\section{Introduction}
\label{sec:intro}

Multimodal generative models must reconcile data with very different structures. Images and audio
are commonly modelled through continuous diffusion, whereas language and protein sequences are
usually generated autoregressively or with discrete corruption processes
\citep{ho2020denoising,song2021score,austin2021structured,sahoo2024simple,wang2024dplm}. Existing
systems bridge these choices through modality-specific encoders, latent spaces, objectives or
prediction interfaces \citep{bao2023unidiffuser,tang2023codi,campbell2024multiflow,wang2025dplm2}.
These designs are effective, but each new modality pairing introduces another interface to design
and calibrate.

We ask whether a single generative interface can instead operate directly on discrete codes from
different modalities. Our starting point is CoBit \citep{method:cobit}, which writes discrete tokens
as binary codes and applies continuous Gaussian diffusion to their bits. We map each modality to
integer codes with a frozen tokenizer, place the codes in a shared binary vocabulary, and arrange
them in a fixed sequence template. Inside the denoiser, every position uses the same input
projection, transformer backbone, bitwise prediction head and loss. A partial-denoising objective
then trains one checkpoint to perform joint generation and every conditional direction included in
training: clean positions specify the conditioning signal, and corrupted positions specify what to
generate. The architecture contains no modality- or segment-specific parameters. We retain CoBit's
probability-space objective, self-conditioning, and data-derived noise scheduling, concentrating
the new method on representation, layout and task selection.

We evaluate this formulation on three complementary modality pairs. On MNIST-Sum, a single
63.3M-parameter denoiser reaches 97.90\% image-to-text accuracy, 99.98\% text-to-image accuracy and
96.85\% joint consistency. \ifdefined\toyimageonly\else On zero-shot MS-COCO evaluation, a
462M-parameter image--text model trained on CC3M obtains FID-30K 18.48 with 20 denoiser evaluations
while supporting captioning and joint generation with the same weights. \fi A 633M-parameter
speech--text model trained without text pretraining supports synthesis, recognition, continuation
and joint generation; in the joint setting, generated speech and text agree to 3.5\% WER and the
speech obtains UTMOS 3.77. A protein model supports folding, inverse folding and joint generation;
80 generated backbones of 100--500 residues reach mean scTM 0.899 after redesign and refolding,
with 68.8\% at or above 0.9. Together, these experiments show that one binary denoising interface
can span raw pixels, learned acoustic codes, text tokens, amino acids and structural codes. They
also identify the remaining constraints imposed by frozen tokenizers, limited text pretraining and
iterative sampling.

Our contributions are a modality-agnostic bitstream interface, a partial-denoising formulation
that selects joint or conditional generation without task-specific modules, and an evaluation of
that formulation across image--text, speech--text and protein sequence--structure generation.
